\documentclass{dlproblemset}
\usepackage{emoji}
\usepackage{booktabs}
\newcommand{\thetav}{\boldsymbol{\theta}}
\DeclareMathOperator*{\argmax}{arg\,max}
\DeclareMathOperator*{\argmin}{arg\,min}

\title{Generative Adversarial Networks}
\subtitle{Lista de Exercícios}

\begin{document}

\paragraph{Instruções:} Responda às questões de múltipla escolha assinalando apenas uma alternativa. Nas dissertativas, apresente o raciocínio passo a passo, incluindo fórmulas e cálculos intermediários quando solicitado. Mantenha a notação dos slides: $g[\mathbf{z},\thetav]$ para o gerador (com latente $\mathbf{z}$ e parâmetros $\thetav$), $f[\mathbf{x},\boldsymbol{\phi}]$ para o discriminador/crítico, $\sigma(\cdot)$ para a sigmóide, $P(\mathbf{x})$ para a distribuição real, $P^{*}(\mathbf{x})$ para a distribuição gerada, $M(\mathbf{x}) = \tfrac{1}{2}(P+P^{*})$ para a mistura, $D^{*}(\mathbf{x}) = \tfrac{P(\mathbf{x})}{P(\mathbf{x})+P^{*}(\mathbf{x})}$ para o discriminador ótimo, $D_{JS}$ para a divergência de Jensen--Shannon e $D_{w}$ para a distância de Wasserstein. Nas questões sobre Wasserstein, as distribuições comparadas são denotadas $P$ e $Q$, e o plano de transporte é $T$, como nos slides. Questões marcadas com \emoji{skull-and-crossbones} são consideradas difíceis.

\subsection*{Questões de Múltipla Escolha.}

\begin{enumerate}[I)]

% ------------------------------------------------------------------
% Q1 — Easy (compute optimal discriminator at a point)
% ------------------------------------------------------------------
\item Considere a categoria $a$, com $P(a) = 0.40$ e $P^{*}(a) = 0.10$. Qual é o valor do discriminador ótimo $D^{*}(a)$ nesse ponto?
\begin{enumerate}[(a)]
    \item $0.80$.
    \item $0.20$.
    \item $0.40$.
    \item $0.50$.
\end{enumerate}
% R: a

% ------------------------------------------------------------------
% Q2 — Easy (BCE contribution for one real sample)
% ------------------------------------------------------------------
\item Para uma única amostra \textbf{real} ($y = 1$), o discriminador devolve $\sigma(f[\mathbf{x},\boldsymbol{\phi}]) = 0.20$. Qual é a contribuição dessa amostra à \emph{loss} de entropia cruzada binária do discriminador?
\begin{enumerate}[(a)]
    \item $0.200$.
    \item $\approx 0.223$.
    \item $\approx 1.609$.
    \item $0.800$.
\end{enumerate}
% R: c

% ------------------------------------------------------------------
% Q3 — Medium (NS vs original gen-loss gradient at d = -3)
% ------------------------------------------------------------------
\item Seja $d = f[g[\mathbf{z},\thetav],\boldsymbol{\phi}]$ o logit do discriminador na amostra gerada. No ponto $d = -3$, quais são, respectivamente, as magnitudes dos gradientes da \emph{loss} original do gerador $L_{\text{orig}} = \ln[1-\sigma(d)]$ e da \emph{loss} \textit{non-saturating} $L_{\text{NS}} = -\ln[\sigma(d)]$ em relação a $d$? Use $\sigma(-3) \approx 0.047$.
\begin{enumerate}[(a)]
    \item $0.953$ e $0.047$.
    \item $0.047$ e $0.953$.
    \item $0.500$ e $0.500$.
    \item $0.047$ e $0.047$.
\end{enumerate}
% R: b

% ------------------------------------------------------------------
% Q4 — Medium (JS divergence over a 2-bin support)
% ------------------------------------------------------------------
\item Sobre o suporte $\{a, b\}$, considere $P = (0.5,\, 0.5)$ e $P^{*} = (0.0,\, 1.0)$. Qual é o valor de $D_{JS}(P^{*}\|P)$ em nats? Use $\ln 2 \approx 0.693$, $\ln(4/3) \approx 0.288$ e $\ln(2/3) \approx -0.405$.
\begin{enumerate}[(a)]
    \item $\approx 0.14$.
    \item $\approx 0.43$.
    \item $\approx 0.69$.
    \item $\approx 0.22$.
\end{enumerate}
% R: d

% ------------------------------------------------------------------
% Q5 — Medium (Wasserstein primal: optimal plan + cost)
% ------------------------------------------------------------------
\item Sobre o suporte $\{0, 1, 2\}$, considere $P = (0.5,\, 0.3,\, 0.2)$ e $Q = (0.2,\, 0.3,\, 0.5)$. Usando o custo $|i - j|$, qual é o valor de $D_{w}(P\|Q)$ obtido pelo plano de transporte ótimo?
\begin{enumerate}[(a)]
    \item $0.0$.
    \item $1.2$.
    \item $0.3$.
    \item $0.6$.
\end{enumerate}
% R: d

% ------------------------------------------------------------------
% Q6 — Medium (DCGAN: replace Tanh output by ReLU — consequence)
% ------------------------------------------------------------------
\item Você implementa uma DCGAN cujas imagens reais estão normalizadas em $[-1, 1]$. Por engano, troca a ativação \emph{Tanh} da última camada do gerador por \emph{ReLU}, mantendo o resto da arquitetura. Qual o efeito mais provável no treinamento?
\begin{enumerate}[(a)]
    \item ReLU produz amostras de qualidade semelhante à \emph{Tanh} desde que haja \emph{BatchNorm} imediatamente antes da última convolução.
    \item O gerador fica restrito a $[0, \infty)$ e nunca produz os valores negativos das imagens reais, e o treinamento degenera.
    \item ReLU acelera o treinamento por eliminar o \emph{vanishing gradient} típico da \emph{Tanh} em saturação.
    \item Como o gerador da DCGAN não usa camadas \emph{fully connected}, a ativação da última convolução não influencia o intervalo da saída.
\end{enumerate}
% R: b

% ------------------------------------------------------------------
% Q7 — Medium (WGAN-GP penalty value)
% ------------------------------------------------------------------
\item Em \textit{WGAN-GP}, o termo de penalidade adicionado à \emph{loss} é $\lambda\,(\|\nabla_{\hat{\mathbf{x}}} f[\hat{\mathbf{x}},\boldsymbol{\phi}]\|_{2} - 1)^{2}$ avaliado em um ponto interpolado $\hat{\mathbf{x}}$. Para $\lambda = 10$ e $\|\nabla_{\hat{\mathbf{x}}} f[\hat{\mathbf{x}},\boldsymbol{\phi}]\|_{2} = 3$, qual é a contribuição desse ponto à \emph{loss}?
\begin{enumerate}[(a)]
    \item $4$.
    \item $90$.
    \item $40$.
    \item $20$.
\end{enumerate}
% R: c

% ------------------------------------------------------------------
% Q8 — Hard (JS in disjoint-support limit)
% ------------------------------------------------------------------
\item \emoji{skull-and-crossbones} Sobre o suporte $\{a, b, c, d\}$, considere $P = (0.5,\, 0.5,\, 0,\, 0)$ e $P^{*} = (0,\, 0,\, 0.4,\, 0.6)$. Qual é o valor de $D_{JS}(P^{*}\|P)$ em nats? Use a convenção $0\ln 0 = 0$.
\begin{enumerate}[(a)]
    \item $\approx 0.35$.
    \item $\approx 1.39$.
    \item $0$.
    \item $\approx 0.69$.
\end{enumerate}
% R: d

% ------------------------------------------------------------------
% Q9 — Medium (trilemma positioning)
% ------------------------------------------------------------------
\item Em uma equipe de pesquisa, alguém descreve um modelo gerador com a seguinte combinação de propriedades: (i) gera amostras nítidas em uma única passagem pela rede, (ii) frequentemente ignora modas inteiras do conjunto real durante o treinamento, (iii) o ajuste de hiperparâmetros é notoriamente instável. A qual família o modelo descrito corresponde, dentro do trilema apresentado nos slides?
\begin{enumerate}[(a)]
    \item GAN.
    \item Modelo difusor (\textit{diffusion model}).
    \item VAE.
    \item \textit{Normalizing flow}.
\end{enumerate}
% R: a

% ------------------------------------------------------------------
% Q10 — Hard (Wasserstein dual: pick the f that achieves D_w)
% ------------------------------------------------------------------
\item \emoji{skull-and-crossbones} Considere $P = (0.6,\, 0.3,\, 0.1)$ e $Q = (0.1,\, 0.3,\, 0.6)$ sobre o suporte $\{0, 1, 2\}$, cuja distância de Wasserstein é $D_{w} = 1$. Qual dos candidatos abaixo para o crítico $f = (f_0, f_1, f_2)$ satisfaz a restrição 1-Lipschitz $|f_{i+1} - f_i| \le 1$ \emph{e} atinge $D_{w}$ no dual $\sum_{i} P(i)\,f_i - \sum_{j} Q(j)\,f_j$?
\begin{enumerate}[(a)]
    \item $f = (0,\, 1,\, 2)$.
    \item $f = (2,\, 1,\, 0)$.
    \item $f = (10,\, 5,\, 0)$.
    \item $f = (1,\, 1,\, 1)$.
\end{enumerate}
% R: b

\end{enumerate}

\subsection*{Gabarito - Objetivas}
\begin{enumerate}[I)]
    \item a
    \item c
    \item b
    \item d
    \item d
    \item b
    \item c
    \item d
    \item a
    \item b
\end{enumerate}

\subsection*{Resoluções - Objetivas}
\color{blue}
\begin{enumerate}[I)]
    \item $D^{*}(a) = \tfrac{P(a)}{P(a) + P^{*}(a)} = \tfrac{0.40}{0.40 + 0.10} = \tfrac{0.40}{0.50} = 0.80$. A (b) corresponde a inverter os papéis no numerador, $\tfrac{P^{*}}{P + P^{*}} = 0.20$. A (c) corresponde a devolver $P(a)$ sem normalizar pela soma. A (d) corresponde ao palpite \emph{chance} de $0.5$, ignorando os dados.

    \item Como o rótulo é $y = 1$ (real), a contribuição é $-\ln[\sigma(f[\mathbf{x},\boldsymbol{\phi}])] = -\ln(0.20) \approx 1.609$ (alternativa c). A (a) reporta $\sigma$ diretamente, sem aplicar o logaritmo. A (b) usa $-\ln(1 - \sigma) = -\ln(0.80) \approx 0.223$, que seria correto se a amostra fosse gerada ($y = 0$). A (d) reporta $1 - \sigma$, também sem logaritmo.

    \item Diferenciando termo a termo: $\partial L_{\text{orig}}/\partial d = -\sigma(d)$, com módulo $\sigma(d) = 0.047$. Para a \textit{non-saturating}, $\partial L_{\text{NS}}/\partial d = -(1 - \sigma(d))$, com módulo $1 - 0.047 = 0.953$ (alternativa b). A (a) troca os dois rótulos. A (c) assume gradiente constante de $0.5$, que vale apenas no equilíbrio $\sigma = \tfrac{1}{2}$. A (d) atribui o mesmo valor a ambas, ignorando que a NS foi construída exatamente para não saturar quando $\sigma \to 0$.

    \item Com $M = (0.25,\, 0.75)$: $D_{KL}(P\|M) = 0.5\ln 2 + 0.5\ln(2/3) \approx 0.347 - 0.203 = 0.144$, e $D_{KL}(P^{*}\|M) = \ln(4/3) \approx 0.288$ (o termo em $a$ vale $0$ pela convenção $0\ln 0 = 0$). Logo $D_{JS} = \tfrac{1}{2}(0.144 + 0.288) \approx 0.22$ nats (alternativa d). A (b) corresponde a somar as duas KL sem o fator $\tfrac{1}{2}$ externo. A (c) corresponde a tratar as distribuições como suporte disjunto e usar o limite $\ln 2$. A (a) corresponde a usar apenas $D_{KL}(P\|M)$.

    \item As diferenças $P - Q = (+0.3,\, 0,\, -0.3)$ indicam que $0.3$ de massa deve sair do índice $0$ e chegar ao índice $2$. O plano ótimo concentra esse movimento em uma única rota, custo $0.3 \cdot |0 - 2| = 0.6$ (alternativa d). A (b) corresponde a contabilizar duas vezes o mesmo transporte (origem e destino). A (c) corresponde a usar distância $1$ em vez de $2$ entre os extremos. A (a) corresponde ao plano \emph{diagonal} ($T_{ij}$ não-nulo apenas em $i = j$), que viola as restrições marginais quando $P \neq Q$.

    \item Imagens reais ocupam $[-1, 1]$; um gerador com \emph{ReLU} na saída só consegue produzir valores em $[0, \infty)$, e em particular nunca em $[-1, 0)$. O discriminador descobre essa pista determinística (qualquer pixel negativo $\Rightarrow$ real) e fica próximo de $\sigma \approx 0$ nas amostras geradas, fazendo o gradiente da \emph{loss} original do gerador colapsar (alternativa b). A (a) ignora a incompatibilidade entre o intervalo de saída e o domínio dos dados. A (c) confunde \emph{vanishing gradient} de saturação interna com o problema de \emph{range} da saída, são fenômenos distintos. A (d) confunde a recomendação da DCGAN, ausência de \emph{fully connected} ocultas, com indiferença em relação à ativação final (que determina o intervalo da saída).

    \item Penalidade $= \lambda\,(\|\nabla f\|_{2} - 1)^{2} = 10 \cdot (3 - 1)^{2} = 10 \cdot 4 = 40$ (alternativa c). A (b) corresponde a $\lambda\|\nabla f\|^{2} = 10 \cdot 9 = 90$, esquecendo o $-1$ dentro do quadrado. A (a) corresponde a $(3 - 1)^{2} = 4$, esquecendo $\lambda$. A (d) corresponde a $\lambda(3 - 1) = 20$, esquecendo de elevar ao quadrado.

    \item Com $M = (0.25,\, 0.25,\, 0.20,\, 0.30)$: nos pontos em que $P > 0$, $\tfrac{P}{M} = 2$; nos pontos em que $P^{*} > 0$, $\tfrac{P^{*}}{M} = 2$. Logo $D_{KL}(P\|M) = (0.5 + 0.5)\ln 2 = \ln 2$ e, analogamente, $D_{KL}(P^{*}\|M) = (0.4 + 0.6)\ln 2 = \ln 2$. Portanto $D_{JS} = \tfrac{1}{2}(\ln 2 + \ln 2) = \ln 2 \approx 0.69$ nats, o teto da Jensen--Shannon (alternativa d). A (b) corresponde a somar as duas KL sem o fator $\tfrac{1}{2}$. A (c) corresponde a confundir suporte disjunto com distribuições iguais. A (a) corresponde a aplicar o fator $\tfrac{1}{2}$ duas vezes.

    \item A combinação ``rápido na amostragem $+$ má cobertura $+$ instável'' é o canto \emph{qualidade $+$ velocidade} do trilema, ocupado por GANs. A (b) descreve difusores, que estão no canto \emph{qualidade $+$ cobertura} (custosos na amostragem). A (c) e a (d) descrevem famílias do canto \emph{velocidade $+$ cobertura}, que costumam ter cobertura boa e estabilidade alta, em troca de qualidade mais limitada.

    \item Verificando cada candidato: $f = (2, 1, 0)$ (alternativa b) é 1-Lipschitz ($|1 - 2| = |0 - 1| = 1$) e produz $(2 \cdot 0.6 + 1 \cdot 0.3 + 0 \cdot 0.1) - (2 \cdot 0.1 + 1 \cdot 0.3 + 0 \cdot 0.6) = 1.5 - 0.5 = 1$, atingindo $D_{w}$. A (a) é 1-Lipschitz mas tem orientação invertida, $0.5 - 1.5 = -1$; o dual exige o máximo, não o mínimo. A (c) atinge soma $5$, porém viola Lipschitz ($|5 - 10| = 5 > 1$), ficando fora da região viável. A (d) é constante: 1-Lipschitz, mas produz soma $0$, não distinguindo $P$ de $Q$.
\end{enumerate}
\color{black}

\newpage

\subsection*{Questões Dissertativas.}

\begin{enumerate}[I)]

% ------------------------------------------------------------------
% D1 — Medium (derive loss = 2 D_JS - 2 ln 2 with optimal D)
% ------------------------------------------------------------------
\item Partindo da \emph{loss} normalizada do discriminador,
\[
  L[\boldsymbol{\phi}] = \mathbb{E}_{\mathbf{x}^{*}}\!\left[\ln[1 - \sigma(f[\mathbf{x}^{*},\boldsymbol{\phi}])]\right]
          + \mathbb{E}_{\mathbf{x}}\!\left[\ln[\sigma(f[\mathbf{x},\boldsymbol{\phi}])]\right],
\]
mostre, em passos algébricos, que com o discriminador ótimo $D^{*}(\mathbf{x}) = \tfrac{P(\mathbf{x})}{P(\mathbf{x}) + P^{*}(\mathbf{x})}$ a \emph{loss} se reduz a
\[
  L[\boldsymbol{\phi}] = 2\,D_{JS}(P^{*}\|P) - 2\ln 2.
\]
Indique explicitamente: (a) onde a hipótese de iguais prioris ($P(\text{real}) = P(\text{gen}) = \tfrac{1}{2}$) entra; (b) onde aparece o ``truque do fator 2''; (c) por que a constante $-2\ln 2$ não afeta o $\argmax$ em relação a $\boldsymbol{\phi}$ nem o $\argmin$ em relação a $\thetav$ (lembre que, escrita sem o sinal de menos, essa é a forma que o discriminador maximiza e o gerador minimiza).

% ------------------------------------------------------------------
% D3 — Hard (non-saturating loss gradient analysis)
% ------------------------------------------------------------------
\item \emoji{skull-and-crossbones} Seja $d = f[g[\mathbf{z},\thetav],\boldsymbol{\phi}]$ o logit do discriminador na amostra gerada. Defina
\[
  L_{\text{orig}}[\thetav] = \ln[1 - \sigma(d)] \qquad\text{e}\qquad L_{\text{NS}}[\thetav] = -\ln[\sigma(d)],
\]
ambas a serem minimizadas pelo gerador.
\begin{enumerate}[(a)]
    \item Derive, mostrando todos os passos, que $\partial L_{\text{orig}}/\partial d = -\sigma(d)$ e $\partial L_{\text{NS}}/\partial d = -(1 - \sigma(d))$. (Lembre que $\sigma'(d) = \sigma(d)(1 - \sigma(d))$.)
    \item Mostre que as duas formulações apontam o gradiente \emph{na mesma direção} (ambas aumentam $d$) e que, no equilíbrio adversarial $\sigma(d) = \tfrac{1}{2}$, as magnitudes dos gradientes coincidem e valem $\tfrac{1}{2}$ em ambos os casos.
    \item Avalie ambas as magnitudes em $d \in \{-6,\, 0,\, +6\}$ e organize em uma pequena tabela. Comente, em uma frase, o caso $d \to -\infty$ (discriminador vencendo o jogo): qual das duas \emph{losses} produz gradiente útil para treinar o gerador?
\end{enumerate}

% ------------------------------------------------------------------
% D4 — Medium (quality vs coverage decomposition, mode collapse)
% ------------------------------------------------------------------
\item A divergência de Jensen--Shannon entre real e gerada decompõe-se como
\[
  D_{JS}(P^{*}\|P) = \underbrace{\tfrac{1}{2}D_{KL}\!\left[P^{*}\,\Big\|\,M\right]}_{\text{qualidade}}
                   + \underbrace{\tfrac{1}{2}D_{KL}\!\left[P\,\Big\|\,M\right]}_{\text{cobertura}},
                   \quad M = \tfrac{1}{2}(P + P^{*}).
\]
\begin{enumerate}[(a)]
    \item Para $P = (0.5,\, 0.3,\, 0.2)$ e $P^{*} = (0.2,\, 0.5,\, 0.3)$ sobre $\{a, b, c\}$, calcule $M$, $D_{KL}(P\|M)$, $D_{KL}(P^{*}\|M)$ e $D_{JS}$. Use $\ln 2 \approx 0.693$ e apresente o resultado em nats com três casas decimais.
    \item Suponha que, em alguma região $R$, $P^{*}(\mathbf{x}) \ll P(\mathbf{x})$ (o gerador ignora uma moda). Mostre que, nessa região, o integrando do termo de \emph{cobertura} é aproximadamente $P(\mathbf{x})\ln 2$, ou seja, fica \textbf{limitado e pouco sensível} a pequenos aumentos de $P^{*}(\mathbf{x})$.
    \item Em duas ou três frases, conecte (b) ao fenômeno de \textit{mode collapse}: por que o gerador prefere ``ganhar qualidade'' nas modas que já cobre a ``ganhar cobertura'' em modas que ignora?
\end{enumerate}

\end{enumerate}

\subsection*{Gabarito}
\color{blue}
\begin{enumerate}[I)]
    \item Partimos de
    \[
      L[\boldsymbol{\phi}] = \mathbb{E}_{\mathbf{x}^{*}}\!\left[\ln[1 - \sigma(f[\mathbf{x}^{*},\boldsymbol{\phi}])]\right]
              + \mathbb{E}_{\mathbf{x}}\!\left[\ln[\sigma(f[\mathbf{x},\boldsymbol{\phi}])]\right]
              = \int P^{*}(\mathbf{x})\ln[1 - \sigma(f)]\,d\mathbf{x} + \int P(\mathbf{x})\ln[\sigma(f)]\,d\mathbf{x}.
    \]
    O discriminador ótimo é obtido pela regra de Bayes assumindo $P(\text{real}) = P(\text{gen}) = \tfrac{1}{2}$ \textbf{(item (a))}, o que cancela os \emph{priors} e produz $\sigma(f[\mathbf{x},\boldsymbol{\phi}]) = \tfrac{P(\mathbf{x})}{P(\mathbf{x}) + P^{*}(\mathbf{x})}$, com complemento $1 - \sigma = \tfrac{P^{*}(\mathbf{x})}{P(\mathbf{x}) + P^{*}(\mathbf{x})}$. Substituindo,
    \[
      L = \int P^{*}(\mathbf{x})\ln\!\left[\tfrac{P^{*}(\mathbf{x})}{P(\mathbf{x}) + P^{*}(\mathbf{x})}\right]d\mathbf{x} + \int P(\mathbf{x})\ln\!\left[\tfrac{P(\mathbf{x})}{P(\mathbf{x}) + P^{*}(\mathbf{x})}\right]d\mathbf{x}.
    \]
    Introduzimos a mistura $M(\mathbf{x}) = \tfrac{1}{2}(P(\mathbf{x}) + P^{*}(\mathbf{x}))$, de modo que $P(\mathbf{x}) + P^{*}(\mathbf{x}) = 2M(\mathbf{x})$. Multiplicando e dividindo por $2$ dentro de cada logaritmo \textbf{(item (b))}, cada razão se reescreve como $\tfrac{2P^{*}(\mathbf{x})}{P(\mathbf{x}) + P^{*}(\mathbf{x})} = \tfrac{P^{*}(\mathbf{x})}{M(\mathbf{x})}$ (e analogamente para $P(\mathbf{x})$):
    \[
      L = \int P^{*}(\mathbf{x})\!\left(\ln\!\left[\tfrac{P^{*}(\mathbf{x})}{M(\mathbf{x})}\right] - \ln 2\right)d\mathbf{x}
        + \int P(\mathbf{x})\!\left(\ln\!\left[\tfrac{P(\mathbf{x})}{M(\mathbf{x})}\right] - \ln 2\right)d\mathbf{x}.
    \]
    Os termos $-\ln 2$ saem das integrais e somam $-2\ln 2$, pois $\int P^{*}(\mathbf{x})\,d\mathbf{x} = \int P(\mathbf{x})\,d\mathbf{x} = 1$. As integrais restantes são exatamente $D_{KL}(P^{*}\|M)$ e $D_{KL}(P\|M)$. Logo
    \[
      L = D_{KL}(P^{*}\|M) + D_{KL}(P\|M) - 2\ln 2 = 2\,D_{JS}(P^{*}\|P) - 2\ln 2.
    \]
    \textbf{Item (c)}: $-2\ln 2$ é constante em $\boldsymbol{\phi}$ e em $\thetav$. Como essa forma da \emph{loss} (sem o sinal de menos) é a log-verossimilhança que o discriminador \emph{maximiza} e que o gerador \emph{minimiza} (minimizar $L$ equivale a minimizar $D_{JS}$), a constante não altera $\argmax_{\boldsymbol{\phi}}$ nem $\argmin_{\thetav}$.

    \item \textbf{(a)} Usando $\sigma'(d) = \sigma(d)(1 - \sigma(d))$:
    \[
      \tfrac{\partial L_{\text{orig}}}{\partial d}
      = \tfrac{\partial}{\partial d}\ln[1 - \sigma(d)]
      = \tfrac{-\sigma'(d)}{1 - \sigma(d)}
      = \tfrac{-\sigma(1 - \sigma)}{1 - \sigma}
      = -\sigma(d).
    \]
    \[
      \tfrac{\partial L_{\text{NS}}}{\partial d}
      = -\tfrac{\partial}{\partial d}\ln[\sigma(d)]
      = -\tfrac{\sigma'(d)}{\sigma(d)}
      = -\tfrac{\sigma(1 - \sigma)}{\sigma}
      = -(1 - \sigma(d)).
    \]

    \textbf{(b)} Ambas as derivadas são \emph{negativas} para $\sigma(d) \in (0, 1)$, portanto a descida de gradiente em $d$ aumenta $d$ em ambos os casos: as duas \emph{losses} empurram o gerador na mesma direção, mudando apenas a intensidade. No equilíbrio adversarial $\sigma(d) = \tfrac{1}{2}$ (ponto em que o discriminador não distingue real de gerado), as duas magnitudes coincidem: $|\partial L_{\text{orig}}/\partial d| = \sigma(d) = \tfrac{1}{2}$ e $|\partial L_{\text{NS}}/\partial d| = 1 - \sigma(d) = \tfrac{1}{2}$.

    \textbf{(c)}
    \begin{center}
    \begin{tabular}{cccc}
      \toprule
      $d$ & $\sigma(d)$ & $|\partial L_{\text{orig}}/\partial d| = \sigma(d)$ & $|\partial L_{\text{NS}}/\partial d| = 1 - \sigma(d)$ \\
      \midrule
      $-6$ & $\approx 0.0025$ & $\approx 0.0025$ & $\approx 0.9975$ \\
      $\;\;0$ & $0.5000$ & $0.5000$ & $0.5000$ \\
      $+6$ & $\approx 0.9975$ & $\approx 0.9975$ & $\approx 0.0025$ \\
      \bottomrule
    \end{tabular}
    \end{center}
    Quando $d \to -\infty$ (discriminador vencendo, $\sigma \to 0$), $|\partial L_{\text{orig}}/\partial d| \to 0$ e a \emph{loss} original satura, ao passo que $|\partial L_{\text{NS}}/\partial d| \to 1$ e o gerador continua recebendo gradiente útil; por isso a NS é o \emph{default} na prática.

    \item \textbf{(a)} $M = \tfrac{1}{2}(P + P^{*}) = (0.35,\, 0.40,\, 0.25)$.
    \[
      D_{KL}(P\|M) = 0.5\ln\tfrac{0.5}{0.35} + 0.3\ln\tfrac{0.3}{0.40} + 0.2\ln\tfrac{0.2}{0.25}
                   \approx 0.178 - 0.086 - 0.045 = 0.047.
    \]
    \[
      D_{KL}(P^{*}\|M) = 0.2\ln\tfrac{0.2}{0.35} + 0.5\ln\tfrac{0.5}{0.40} + 0.3\ln\tfrac{0.3}{0.25}
                       \approx -0.112 + 0.112 + 0.055 = 0.054.
    \]
    \[
      D_{JS}(P^{*}\|P) = \tfrac{1}{2}(0.047 + 0.054) \approx 0.051\ \text{nats}.
    \]

    \textbf{(b)} Em região $R$ onde $P^{*}(\mathbf{x}) \ll P(\mathbf{x})$, temos $M(\mathbf{x}) = \tfrac{P(\mathbf{x}) + P^{*}(\mathbf{x})}{2} \approx \tfrac{P(\mathbf{x})}{2}$. Logo $\tfrac{P(\mathbf{x})}{M(\mathbf{x})} \approx 2$, e o integrando da cobertura fica
    \[
      P(\mathbf{x})\ln\!\tfrac{P(\mathbf{x})}{M(\mathbf{x})} \approx P(\mathbf{x})\ln 2,
    \]
    \textbf{limitado} (cresce apenas com $P(\mathbf{x})$, não com a discrepância em si) e \textbf{pouco sensível} a pequenas perturbações em $P^{*}(\mathbf{x})$ (a derivada em $P^{*}(\mathbf{x})$ envolve $\tfrac{1}{P(\mathbf{x}) + P^{*}(\mathbf{x})}$, que é pequena quando $P(\mathbf{x})$ domina).

    \textbf{(c)} O termo de qualidade ($D_{KL}(P^{*}\|M)$) explode se o gerador colocar massa em pontos onde $P$ é baixa: o gradiente pune \textbf{fortemente} amostras implausíveis. O termo de cobertura, pelo argumento de (b), é \textbf{quase insensível} a aumentos pequenos de $P^{*}(\mathbf{x})$ em modas ignoradas. Diante dessa assimetria, o caminho de menor custo durante o treinamento é \emph{melhorar a qualidade nas modas já cobertas} (penalidade alta e gradiente forte) em vez de \emph{cobrir novas modas} (penalidade baixa e gradiente fraco), produzindo \textit{mode collapse}.
\end{enumerate}

\end{document}
